\section{Discussion}
\label{sec:discussion}

\subsection{Relation to Vol.~1 and Vol.~2}

Vol.~1 established time and space as emergent from the causal log $C_{ij}$.
Vol.~2 showed that particle physics structure emerges from triangular closure.
The present volume completes this picture by deriving gravitational dynamics
from closure density.

The key unifying element is the closure energy functional.
In Vol.~1, it governed the emergence of metric structure.
In Vol.~2, its local minima were identified with particle types.
In Vol.~3, its coarse-grained form determines the gravitational field.
All three volumes therefore share a common geometric foundation.

\subsection{Scope and Limitations}

The interpretations advanced in this volume are structural and qualitative.
Several points of caution are warranted.

First, the identification of $G_{\mu\nu} = 8\pi G_{\text{eff}} L_{\mu\nu}$
with the Einstein equation is structural.
No derivation is given of the precise numerical coefficient $8\pi$,
nor of the relation between $G_{\text{eff}}$ and observed gravitational coupling.
A rigorous derivation would require a more developed formalism.

Second, the Schwarzschild metric is recovered only in the weak-gravity limit.
The full nonlinear solution is asserted to emerge from the exact fixed-point
equation, but this has not been demonstrated explicitly.

Third, the cosmological equation~\eqref{eq:VED_cosmo} is structurally
analogous to the Friedmann equation, but the precise correspondence
between $\rho$ and energy density, and between $\gamma$ and the equation
of state, has not been established.

Fourth, the claim that the black hole interior corresponds to a
$\sigma > 0$ phase, while motivated by the prohibition of perfect closure,
remains a qualitative assertion.
A quantitative treatment would require extending the fixed-point formalism
into the non-equilibrium regime.

\subsection{Relation to General Relativity}

The present framework does not compete with general relativity.
GR is accurate in the classical closure phase ($\sigma \approx 0$),
and the present work is consistent with this.
The distinction arises in other phases:
at the black hole interior, at the early universe,
and at any regime where the fixed-point structure breaks down.

\subsection{Path Forward}

Several directions suggest themselves for future development.

The most immediate is a more rigorous treatment of the continuum limit:
deriving the closure tensor $L_{\mu\nu}$ explicitly from the discrete network,
and showing that the linearized fixed-point equation takes the form
$G_{\mu\nu} = 8\pi G_{\text{eff}} L_{\mu\nu}$ without additional assumptions.

A second direction is the cosmological model.
The logistic form of $\sigma(\rho)$ is motivated by self-referential selection,
but its precise form should be derivable from the fixed-point structure.
Connecting the fixed-point density $\rho_*$ to observed cosmological parameters
would constitute a non-trivial test of the framework.

A third direction concerns the non-equilibrium phase.
The interior dynamics of $\sigma > 0$ regions — including the black hole interior
and the early universe — require a treatment beyond the classical fixed-point
approximation. This may connect to information-theoretic questions about
black hole evaporation and the origin of entropy.
